% SPDX-FileCopyrightText: © 2026 Idiap Research Institute <contact@idiap.ch>
%
% SPDX-FileContributor: Andre Freitas <andre.freitas@idiap.ch>
% Neuro-symbolic AI Group
%
% SPDX-FileContributor: Danilo Gusicuma <danilo.gusicuma@idiap.ch>
% Neuro-symbolic AI Group
%
% SPDX-License-Identifier: MIT
% Options for packages loaded elsewhere
\PassOptionsToPackage{unicode}{hyperref}
\PassOptionsToPackage{hyphens}{url}
\documentclass[
]{article}
\usepackage{xcolor}
\usepackage{amsmath,amssymb}
\setcounter{secnumdepth}{-\maxdimen} % remove section numbering
\usepackage{iftex}
\ifPDFTeX
  \usepackage[T1]{fontenc}
  \usepackage[utf8]{inputenc}
  \usepackage{textcomp} % provide euro and other symbols
\else % if luatex or xetex
  \usepackage{unicode-math} % this also loads fontspec
  \defaultfontfeatures{Scale=MatchLowercase}
  \defaultfontfeatures[\rmfamily]{Ligatures=TeX,Scale=1}
\fi
\usepackage{lmodern}
\ifPDFTeX\else
  % xetex/luatex font selection
\fi
% Use upquote if available, for straight quotes in verbatim environments
\IfFileExists{upquote.sty}{\usepackage{upquote}}{}
\IfFileExists{microtype.sty}{% use microtype if available
  \usepackage[]{microtype}
  \UseMicrotypeSet[protrusion]{basicmath} % disable protrusion for tt fonts
}{}
\makeatletter
\@ifundefined{KOMAClassName}{% if non-KOMA class
  \IfFileExists{parskip.sty}{%
    \usepackage{parskip}
  }{% else
    \setlength{\parindent}{0pt}
    \setlength{\parskip}{6pt plus 2pt minus 1pt}}
}{% if KOMA class
  \KOMAoptions{parskip=half}}
\makeatother
\setlength{\emergencystretch}{3em} % prevent overfull lines
\providecommand{\tightlist}{%
  \setlength{\itemsep}{0pt}\setlength{\parskip}{0pt}}
\usepackage{bookmark}
\IfFileExists{xurl.sty}{\usepackage{xurl}}{} % add URL line breaks if available
\urlstyle{same}
\hypersetup{
  pdftitle={Yang Mills Existence And Mass Gap.hypotheses Evaluation Explained},
  hidelinks,
  pdfcreator={LaTeX via pandoc}}

\title{Yang Mills Existence And Mass Gap.hypotheses Evaluation
Explained}
\author{}
\date{}

\begin{document}
\maketitle

\section{Yang-\/-Mills Hypothesis Evaluation
Explained}\label{yang--mills-hypothesis-evaluation-explained}

\subsection{Purpose}\label{purpose}

This document explains how to read the hypothesis evaluation report in
\href{/Users/user/files/abducer/abducer/phyp/yang-mills-existence-and-mass-gap.hypotheses-evaluated.md}{yang-mills-existence-and-mass-gap.hypotheses-evaluated.md}.
The report is not a proof checker. It is a triage layer over proof
hypotheses generated from the Yang-\/-Mills knowledge base. Its job is
to decide which hypotheses are currently well-grounded in the KB, which
ones remain plausible but underdetermined, and which ones should be
deprioritized or refuted.

The central point is this:

\begin{itemize}
\tightlist
\item
  \texttt{supported} means the framework currently sees enough
  structural and coherence-based evidence to keep the hypothesis in the
  active frontier
\item
  \texttt{inconclusive} means the hypothesis is still live, but the
  present KB does not yet justify promoting it
\item
  \texttt{refuted} would mean the framework found sufficiently strong
  structural conflict or explicit negative evidence
\end{itemize}

In this run there are no \texttt{refuted} hypotheses. That matters. The
system is not saying that most proposals are bad. It is saying that most
proposals are not yet sharp enough, evidenced enough, or deductively
connected enough to the current Yang-\/-Mills KB.

\subsection{What Was Evaluated}\label{what-was-evaluated}

\begin{itemize}
\tightlist
\item
  Total hypotheses: \texttt{44}
\item
  Direct bridge hypotheses: \texttt{001} to \texttt{020}
\item
  Intervention-focus hypotheses: \texttt{021} to \texttt{032}
\item
  Higher-order conceptual hypotheses: \texttt{033} to \texttt{044}
\end{itemize}

These are three different inferential roles.

\begin{itemize}
\tightlist
\item
  \texttt{bridge-link}: a direct proposed connection from one statement
  to another
\item
  \texttt{intervention-focus}: a claim that a particular statement is a
  high-value place to concentrate reasoning effort
\item
  \texttt{representation-shift} / \texttt{intermediate-invariant} /
  \texttt{definition-introduction}: higher-order conceptual moves that
  reorganize the proof search space rather than asserting a single
  direct theorem jump
\end{itemize}

This distinction is essential for interpreting the result. The strongest
results in this run are mostly not direct theorem-to-theorem bridges.
They are conceptual restructurings.

\subsection{How The Evaluator Decides}\label{how-the-evaluator-decides}

The evaluator combines four layers of evidence.

\subsubsection{1. Structural support}\label{1-structural-support}

For bridge and conceptual hypotheses, the evaluator looks at:

\begin{itemize}
\tightlist
\item
  typed similarity between source and target statement structure
\item
  symbol-role alignment
\item
  theorem-family similarity
\item
  connector analysis score
\item
  whether the target already depends on the source
\item
  whether the proposal loops back on itself
\item
  whether there are contradiction markers
\end{itemize}

This is implemented in
\href{/Users/user/files/abducer/abducer/skills/common/hypothesis_eval.py:755}{hypothesis\_eval.py}
and
\href{/Users/user/files/abducer/abducer/skills/common/hypothesis_eval.py:969}{hypothesis\_eval.py}.

\subsubsection{2. Inferential commitment
grounding}\label{2-inferential-commitment-grounding}

Each non-intervention hypothesis can carry four explicit commitments:

\begin{itemize}
\tightlist
\item
  antecedent schema
\item
  target consequence schema
\item
  preserved invariant
\item
  barrier scope
\end{itemize}

The evaluator asks whether those commitments are actually visible in the
structured content of the source and target statements. This is why the
report contains alignment signals such as:

\begin{itemize}
\tightlist
\item
  \texttt{antecedent\_alignment}
\item
  \texttt{consequence\_alignment}
\item
  \texttt{invariant\_alignment}
\item
  \texttt{barrier\_alignment}
\end{itemize}

This matters because the framework is trying to reason mathematically,
not just lexically. A hypothesis is stronger when it states a concrete
inferential interface and that interface is visible in the KB statements
themselves.

\subsubsection{3. External evidence
adjustments}\label{3-external-evidence-adjustments}

Optional deductive, abductive, contradiction, or experimental signals
can adjust support or refutation. The logic is in
\href{/Users/user/files/abducer/abducer/skills/common/hypothesis_eval.py:620}{hypothesis\_eval.py}.

For example:

\begin{itemize}
\tightlist
\item
  a deductive entailment can increase support
\item
  a failed derivation can increase refutation
\item
  an explicit contradiction can strongly increase refutation
\end{itemize}

In this Yang-\/-Mills run, the results are driven mainly by structural
and coherence signals, not by hard deductive discharge.

\subsubsection{4. Coherence adjustments}\label{4-coherence-adjustments}

Each hypothesis also receives a local coherence readout against the
current theory graph. This is why the report lists:

\begin{itemize}
\tightlist
\item
  \texttt{Coherence\ signals:\ 44}
\item
  \texttt{Global\ coherence\ baseline:\ 0.320043}
\end{itemize}

The baseline is the current global coherence score of the KB before
accepting any specific hypothesis. Each hypothesis is then evaluated
against that background. This does not yet mean the system has found a
globally optimized theory. It means each hypothesis is judged partly by
whether it improves or destabilizes the current web of claims.

\subsection{How Verdicts Are Assigned}\label{how-verdicts-are-assigned}

The decision rule is explicit in
\href{/Users/user/files/abducer/abducer/skills/common/hypothesis_eval.py:662}{hypothesis\_eval.py}.

\begin{itemize}
\tightlist
\item
  \texttt{supported} requires:

  \begin{itemize}
  \tightlist
  \item
    \texttt{support\ \textgreater{}=\ 0.62}
  \item
    and \texttt{support\ -\ refute\ \textgreater{}=\ 0.08}
  \end{itemize}
\item
  \texttt{refuted} would require:

  \begin{itemize}
  \tightlist
  \item
    \texttt{refute\ \textgreater{}=\ 0.62}
  \item
    and \texttt{refute\ -\ support\ \textgreater{}=\ 0.08}
  \end{itemize}
\item
  otherwise the result is \texttt{inconclusive}
\end{itemize}

For this run:

\begin{itemize}
\tightlist
\item
  support threshold = \texttt{0.62}
\item
  refute threshold = \texttt{0.62}
\item
  margin = \texttt{0.08}
\item
  force resolution = \texttt{False}
\end{itemize}

\texttt{force\ resolution\ =\ False} is important. It means the
evaluator is allowed to leave a hypothesis unresolved. It does not have
to pretend every hypothesis is either accepted or rejected.

\subsection{What Confidence Means}\label{what-confidence-means}

\texttt{confidence} is not the same thing as \texttt{support}.

It is computed from:

\begin{itemize}
\tightlist
\item
  the stronger of the support/refute scores
\item
  multiplied by a factor that rewards separation between them
\end{itemize}

The formula is in
\href{/Users/user/files/abducer/abducer/skills/common/hypothesis_eval.py:854}{hypothesis\_eval.py},
\href{/Users/user/files/abducer/abducer/skills/common/hypothesis_eval.py:1019}{hypothesis\_eval.py},
and
\href{/Users/user/files/abducer/abducer/skills/common/hypothesis_eval.py:1170}{hypothesis\_eval.py}.

Interpretation:

\begin{itemize}
\tightlist
\item
  high \texttt{support} with low \texttt{refute} gives higher confidence
\item
  moderate \texttt{support} with moderate \texttt{refute} can still be
  \texttt{supported}, but with lower confidence
\end{itemize}

That is why \texttt{phyp.yang-mills-existence-and-mass-gap.040} is
\texttt{supported} even though its confidence is only \texttt{0.4823}:
its support clears threshold and its margin over refute is large enough,
but the absolute strength is still moderate.

\subsection{What the Summary Means}\label{what-the-summary-means}

\subsubsection{Stats}\label{stats}

\begin{itemize}
\tightlist
\item
  \texttt{44} hypotheses were evaluated
\item
  \texttt{44} coherence signals means every hypothesis had a coherence
  row
\item
  global coherence baseline \texttt{0.320043} means the current KB is
  only moderately integrated; the framework is still working inside a
  partially coherent theory graph
\end{itemize}

\subsubsection{Verdict counts}\label{verdict-counts}

\begin{itemize}
\tightlist
\item
  \texttt{12\ supported}
\item
  \texttt{32\ inconclusive}
\item
  \texttt{0\ refuted}
\end{itemize}

This pattern has a clear interpretation:

\begin{itemize}
\tightlist
\item
  the framework found several promising conceptual reorganizations
\item
  it did not yet find enough grounded direct bridges to turn those
  reorganizations into a strong proof path
\item
  it also did not find decisive negative evidence against most of the
  frontier
\end{itemize}

So the current Yang-\/-Mills state is exploratory but not collapsed.

\subsection{How To Read The Three Hypothesis
Classes}\label{how-to-read-the-three-hypothesis-classes}

\subsection{\texorpdfstring{Intervention-Focus Hypotheses (\texttt{021}
to
\texttt{032})}{Intervention-Focus Hypotheses (021 to 032)}}\label{intervention-focus-hypotheses-021-to-032}

These do not propose a theorem bridge. They propose where attention
should go next.

Examples:

\begin{itemize}
\tightlist
\item
  \texttt{021}: focus on \texttt{st.def.os-euclidean-qft}
\item
  \texttt{022}: focus on \texttt{st.def.mass-gap}
\item
  \texttt{023}: focus on
  \texttt{st.thm.euclidean-correlator-decay-as-a-mass-gap-signatu}
\end{itemize}

Their alignment signals are all zero in the mathematical readout. That
is expected. Intervention hypotheses do not carry explicit antecedent,
consequence, invariant, or barrier commitments in the same way bridge
hypotheses do. They are scored mostly from:

\begin{itemize}
\tightlist
\item
  inherited priority
\item
  cross-component potential
\item
  downstream impact
\item
  betweenness and dependency pressure
\end{itemize}

This logic is in
\href{/Users/user/files/abducer/abducer/skills/common/hypothesis_eval.py:1079}{hypothesis\_eval.py}.

Interpretation of their current outcome:

\begin{itemize}
\tightlist
\item
  they are plausible focal points
\item
  they are not yet strong enough to be promoted as decisive intervention
  commitments
\item
  they should guide the next reasoning round rather than be treated as
  mathematical conclusions
\end{itemize}

\subsection{\texorpdfstring{Direct Bridge Hypotheses (\texttt{001} to
\texttt{020})}{Direct Bridge Hypotheses (001 to 020)}}\label{direct-bridge-hypotheses-001-to-020}

These are the most literal bridging proposals. They connect one KB
statement to another and state an intended invariant.

Examples:

\begin{itemize}
\tightlist
\item
  \texttt{001}: from \texttt{st.def.os-euclidean-qft} toward OS
  reconstruction and Wightman equivalence, preserving
  \texttt{OS-to-Wightman\ mass-gap\ transfer\ in\ Hamiltonian-mass-gap}
\item
  \texttt{003}: from \texttt{st.def.os-euclidean-qft} toward Euclidean
  correlator decay as a mass-gap signature, preserving
  \texttt{correlator-decay\ mass-gap\ witness\ in\ Schwinger-function-decay}
\item
  \texttt{006}: from \texttt{st.def.mass-gap}-adjacent transfer-matrix
  structure toward correlator decay, preserving
  \texttt{transfer-matrix\ spectral\ witness\ in\ reflection-positive-transfer-matrix}
\end{itemize}

These hypotheses are mostly \texttt{inconclusive}.

That does not mean they are wrong. It means the framework sees real
structure but not enough decisive grounding.

Typical reasons:

\begin{itemize}
\tightlist
\item
  good invariant alignment but weak consequence alignment
\item
  good semantic overlap but barrier scope not yet grounded
\item
  connector score present but not strong enough to clear support
  threshold
\item
  bridge is mathematically plausible but still too under-specified
  relative to the current KB
\end{itemize}

Example:

\begin{itemize}
\tightlist
\item
  \texttt{001} has:

  \begin{itemize}
  \tightlist
  \item
    antecedent alignment \texttt{0.3869}
  \item
    invariant alignment \texttt{0.5114}
  \item
    consequence alignment \texttt{0.0000}
  \item
    support \texttt{0.4709}
  \item
    refute \texttt{0.0540}
  \end{itemize}
\end{itemize}

This says:

\begin{itemize}
\tightlist
\item
  the source really does look like an OS/reconstruction antecedent
\item
  the proposed invariant is visible across source and target
\item
  but the target consequence \texttt{positive-self-adjoint-Hamiltonian}
  is not yet grounded strongly enough in the target statement structure
\end{itemize}

So the bridge is mathematically suggestive but not yet inferentially
tight.

\subsection{\texorpdfstring{Higher-Order Conceptual Hypotheses
(\texttt{033} to
\texttt{044})}{Higher-Order Conceptual Hypotheses (033 to 044)}}\label{higher-order-conceptual-hypotheses-033-to-044}

These are the strongest part of the current run. All \texttt{12}
supported hypotheses are in this class.

This is the central result of the report.

The framework is saying that, for this Yang-\/-Mills KB, the best
current progress is not a direct proof bridge. It is a conceptual
reframing of the proof search space.

The supported hypotheses split into three families.

\subsection{Supported Family A: Representation
Shift}\label{supported-family-a-representation-shift}

Supported hypotheses:

\begin{itemize}
\tightlist
\item
  \texttt{033}
\item
  \texttt{036}
\item
  \texttt{037}
\item
  \texttt{041}
\item
  \texttt{043}
\end{itemize}

These propose that the right reasoning move is to shift representation
between:

\begin{itemize}
\tightlist
\item
  OS Euclidean data
\item
  Wightman axiomatic structure
\item
  Hamiltonian/spectral formulations
\item
  constructive QFT examples
\item
  rigorous renormalization group programs
\end{itemize}

Interpretation:

\begin{itemize}
\tightlist
\item
  the framework sees the Euclidean-to-Wightman/Hamiltonian interface as
  a structurally coherent route
\item
  it prefers proof search that moves through representation changes
  instead of trying to jump directly from geometric or gauge-theoretic
  statements to the Clay conclusion
\end{itemize}

The strongest single supported hypothesis is:

\begin{itemize}
\tightlist
\item
  \texttt{041}

  \begin{itemize}
  \tightlist
  \item
    support \texttt{0.7829}
  \item
    refute \texttt{0.0216}
  \item
    confidence \texttt{0.6988}
  \end{itemize}
\end{itemize}

\texttt{041} is a representation shift from
\texttt{st.def.os-euclidean-qft} to the more explicit OS definition
statement. This is significant because the system is effectively saying:

\begin{itemize}
\tightlist
\item
  before pushing further, stabilize the exact OS formulation
\item
  then use that stabilization as the anchor for the mass-gap route
\end{itemize}

This is not yet a proof of a mass gap. It is a diagnosis that
definitional precision at the OS level is one of the strongest currently
supported moves.

\subsection{Supported Family B: Intermediate
Invariant}\label{supported-family-b-intermediate-invariant}

Supported hypotheses:

\begin{itemize}
\tightlist
\item
  \texttt{034}
\item
  \texttt{035}
\item
  \texttt{038}
\item
  \texttt{044}
\end{itemize}

These propose that the proof search should be mediated by explicitly
preserved middle objects, especially:

\begin{itemize}
\tightlist
\item
  correlator-decay mass-gap witnesses
\item
  transfer-matrix spectral witnesses
\item
  mass-gap witness transfer in compact simple gauge group settings
\end{itemize}

Interpretation:

\begin{itemize}
\tightlist
\item
  the framework does not trust a raw statement-to-statement jump
\item
  it wants an invariant that survives translation between formulations
\end{itemize}

This is mathematically appropriate. In hard existence-and-gap problems,
a preserved witness or invariant is usually what makes a bridge usable.

The two most important invariant families here are:

\begin{itemize}
\tightlist
\item
  \texttt{correlator-decay\ -\textgreater{}\ gap}
\item
  \texttt{transfer-matrix\ -\textgreater{}\ spectral\ gap}
\end{itemize}

These are supported because the Yang-\/-Mills KB already contains the
right structural ingredients:

\begin{itemize}
\tightlist
\item
  OS positivity
\item
  reflection positivity
\item
  transfer-matrix language
\item
  exponential decay language
\item
  Hamiltonian gap language
\end{itemize}

So the evaluator is identifying a meaningful inferential spine:

\begin{itemize}
\tightlist
\item
  prove or control the right decay witness
\item
  transfer that witness through reflection positivity / transfer matrix
  structure
\item
  interpret it as a Hamiltonian gap signal
\end{itemize}

\subsection{Supported Family C: Definition
Introduction}\label{supported-family-c-definition-introduction}

Supported hypotheses:

\begin{itemize}
\tightlist
\item
  \texttt{039}
\item
  \texttt{040}
\item
  \texttt{042}
\end{itemize}

These are geometric/gauge-theoretic conceptual moves from the compact
simple gauge group definition toward:

\begin{itemize}
\tightlist
\item
  ADHM instanton classification
\item
  Uhlenbeck gauge fixing and compactness
\item
  BPST instanton existence
\end{itemize}

Interpretation:

\begin{itemize}
\tightlist
\item
  the system sees these as legitimate concept-building moves
\item
  but it treats them as preparatory structure, not as direct mass-gap
  resolution
\end{itemize}

This is the right reading. These statements provide geometric control,
local compactness, or model constructions. They matter because they
enrich the ontology of Yang-\/-Mills objects. But they are not yet by
themselves a route to the positive mass gap in four-dimensional quantum
Yang-\/-Mills.

So these are supported as concept-building steps, not as decisive
closure steps.

\subsection{What The Report Is Really Saying About
Yang-\/-Mills}\label{what-the-report-is-really-saying-about-yang--mills}

The report points to a fairly specific mathematical diagnosis.

\subsubsection{1. The current KB supports conceptual transport better
than direct proof
transport}\label{1-the-current-kb-supports-conceptual-transport-better-than-direct-proof-transport}

The evaluator supports:

\begin{itemize}
\tightlist
\item
  representation shifts
\item
  invariant introduction
\item
  definition sharpening
\end{itemize}

It does not yet support most direct bridge links.

This means the KB is rich enough to describe the right proof ecology,
but not yet rich enough to discharge the main proof obligations.

\subsubsection{2. The most promising mass-gap route currently runs
through the Euclidean/OS/Wightman/spectral
interface}\label{2-the-most-promising-mass-gap-route-currently-runs-through-the-euclideanoswightmanspectral-interface}

The strongest supported cluster is centered on:

\begin{itemize}
\tightlist
\item
  OS Euclidean QFT
\item
  OS reconstruction
\item
  Wightman realization
\item
  transfer-matrix and reflection positivity
\item
  correlator decay as a signature of spectral gap
\end{itemize}

That is a coherent mathematical route. It says the framework currently
believes the best entry into the Clay problem is not through instantons
alone, nor through gauge-fixing alone, but through the problem of
transporting a gap witness across formulations.

\subsubsection{3. The main blockers remain continuum and ultraviolet
issues}\label{3-the-main-blockers-remain-continuum-and-ultraviolet-issues}

Several of the better bridge proposals are barrier-scoped by:

\begin{itemize}
\tightlist
\item
  \texttt{continuum-limit}
\item
  \texttt{ultraviolet-control}
\item
  \texttt{open-construction-gap}
\end{itemize}

This is exactly what one should expect. The lattice or Euclidean objects
can give partial witness structure, but the unresolved problem is
whether that structure survives the full continuum construction in four
dimensions.

\subsubsection{4. Geometry is present, but not yet integrated into the
mass-gap
mechanism}\label{4-geometry-is-present-but-not-yet-integrated-into-the-mass-gap-mechanism}

The supported geometric hypotheses are mostly:

\begin{itemize}
\tightlist
\item
  ADHM
\item
  BPST instantons
\item
  Uhlenbeck compactness
\item
  Donaldson-style moduli control
\end{itemize}

These are mathematically important, but the report is implicitly saying:

\begin{itemize}
\tightlist
\item
  geometry is currently contributing ontology and structural control
\item
  not yet a precise spectral-gap mechanism
\end{itemize}

So the geometry side of the KB needs better coupling to the mass-gap
witness side.

\subsection{Why So Many Hypotheses Are
Inconclusive}\label{why-so-many-hypotheses-are-inconclusive}

There are three main reasons.

\subsubsection{1. The KB still contains more conceptual adjacency than
proof obligation
data}\label{1-the-kb-still-contains-more-conceptual-adjacency-than-proof-obligation-data}

Many hypotheses are plausible because the statements live in the same
conceptual neighborhood. But plausibility is not enough. The evaluator
wants:

\begin{itemize}
\tightlist
\item
  explicit consequence grounding
\item
  explicit invariant grounding
\item
  visible barrier handling
\end{itemize}

When one of those is missing, support tends to stay in the \texttt{0.43}
to \texttt{0.53} range instead of crossing \texttt{0.62}.

\subsubsection{2. Direct bridges are harder than conceptual
reframings}\label{2-direct-bridges-are-harder-than-conceptual-reframings}

A representation shift only needs to show that one formulation can
organize another in a productive way.

A direct bridge-link has to show more:

\begin{itemize}
\tightlist
\item
  why this exact support statement helps this exact target statement
\item
  what is preserved
\item
  why that preservation survives the stated barriers
\end{itemize}

That is a much harder demand, and the current KB often does not yet
encode enough detail for it.

\subsubsection{3. The Yang-\/-Mills problem is hard in exactly the
places the report
highlights}\label{3-the-yang--mills-problem-is-hard-in-exactly-the-places-the-report-highlights}

The pattern of inconclusiveness is not a defect of the evaluator alone.
It reflects the mathematical problem:

\begin{itemize}
\tightlist
\item
  how to construct the full continuum theory
\item
  how to prove a positive mass gap in the physically relevant setting
\item
  how to connect Euclidean, Hamiltonian, spectral, and constructive
  viewpoints without losing the witness
\end{itemize}

That is precisely where the hypotheses stall.

\subsection{How To Read Individual Fields In The Mathematical
Readout}\label{how-to-read-individual-fields-in-the-mathematical-readout}

\subsubsection{Antecedent schema}\label{antecedent-schema}

This is the type of input condition the hypothesis claims to rely on.

Examples:

\begin{itemize}
\tightlist
\item
  \texttt{reflection-positivity-and-transfer-matrix}
\item
  \texttt{constructive-existence}
\item
  \texttt{unconditional}
\end{itemize}

If antecedent alignment is high, the source statement genuinely looks
like the kind of statement the hypothesis says it needs.

\subsubsection{Target consequence}\label{target-consequence}

This is the type of mathematical outcome the hypothesis claims to move
toward.

Examples:

\begin{itemize}
\tightlist
\item
  \texttt{positive-self-adjoint-Hamiltonian}
\item
  \texttt{rigorous-Yang-Mills-construction}
\item
  \texttt{moduli-space-control}
\end{itemize}

If consequence alignment is low or zero, that usually means the target
statement does not yet explicitly encode that consequence strongly
enough.

\subsubsection{Preserved invariant}\label{preserved-invariant}

This is the most important field in many of the supported hypotheses.

Examples:

\begin{itemize}
\tightlist
\item
  \texttt{OS-to-Wightman\ mass-gap\ transfer\ in\ Hamiltonian-mass-gap}
\item
  \texttt{correlator-decay\ mass-gap\ witness\ in\ Schwinger-function-decay}
\item
  \texttt{transfer-matrix\ spectral\ witness\ in\ reflection-positive-transfer-matrix}
\end{itemize}

This field says what must survive the bridge if the bridge is to have
mathematical value.

\subsubsection{Barrier scope}\label{barrier-scope}

This names what kind of obstruction the move is expected to face.

Examples:

\begin{itemize}
\tightlist
\item
  \texttt{continuum-limit}
\item
  \texttt{ultraviolet-control}
\item
  \texttt{gauge-control}
\item
  \texttt{open-construction-gap\ +\ quantum-measure-construction}
\end{itemize}

A bridge with a realistic barrier scope is often better than a bridge
with no barriers named, because it is more mathematically honest. But it
can also be harder to support if the KB does not yet encode
barrier-handling details.

\subsection{Recommended Reading Of The Supported
Hypotheses}\label{recommended-reading-of-the-supported-hypotheses}

If the goal is to use the current output as a research steering
document, the hypotheses should be read in this order.

\subsubsection{Tier 1: Core mass-gap
interface}\label{tier-1-core-mass-gap-interface}

\begin{itemize}
\tightlist
\item
  \texttt{041}
\item
  \texttt{033}
\item
  \texttt{037}
\item
  \texttt{035}
\item
  \texttt{038}
\item
  \texttt{044}
\end{itemize}

These are the most informative hypotheses for the central problem
because they organize the interface between:

\begin{itemize}
\tightlist
\item
  OS formulation
\item
  Wightman/Hamiltonian realization
\item
  correlator decay
\item
  transfer matrix
\item
  spectral gap witness
\end{itemize}

\subsubsection{Tier 2: Construction-oriented
transfer}\label{tier-2-construction-oriented-transfer}

\begin{itemize}
\tightlist
\item
  \texttt{036}
\item
  \texttt{043}
\end{itemize}

These say that constructive QFT examples and Balaban-style RG analysis
are valuable as representation-shift anchors, but they remain
barrier-loaded and do not yet close the four-dimensional problem.

\subsubsection{Tier 3: Geometric ontology
enrichment}\label{tier-3-geometric-ontology-enrichment}

\begin{itemize}
\tightlist
\item
  \texttt{034}
\item
  \texttt{039}
\item
  \texttt{040}
\item
  \texttt{042}
\end{itemize}

These are useful for concept formation and structural control, but they
are secondary unless they are connected more explicitly to a mass-gap
witness mechanism.

\subsection{What To Do Next}\label{what-to-do-next}

The report itself already suggests the next action:

\begin{itemize}
\tightlist
\item
  \texttt{promote-or-retest-with-deduction} for supported hypotheses
\item
  \texttt{run-another-abduce-deduce-round} for inconclusive hypotheses
\end{itemize}

For the Yang-\/-Mills KB, the right next moves are:

\begin{itemize}
\tightlist
\item
  strengthen the OS/reconstruction/Wightman statements with more
  explicit spectral-gap obligations
\item
  encode more explicit transfer principles from correlator decay to
  Hamiltonian gap
\item
  connect geometric control statements to the mass-gap witness layer
  instead of leaving them as parallel concept clusters
\item
  add sharper continuum-limit and ultraviolet-control statements so
  barrier alignment is no longer mostly implicit
\end{itemize}

\subsection{Bottom Line}\label{bottom-line}

This evaluation does not yet identify a completed Yang-\/-Mills proof
pathway.

It does identify something more precise than a vague frontier:

\begin{itemize}
\tightlist
\item
  the current best-supported discovery path is conceptual and
  representational
\item
  the strongest bridge objects are invariant-like witnesses, not direct
  theorem jumps
\item
  the Euclidean/OS/Wightman/transfer-matrix/correlator-decay chain is
  the main supported inferential spine
\item
  the geometry/moduli/instanton side is useful, but still insufficiently
  coupled to the actual mass-gap mechanism
\end{itemize}

So the correct reading is:

\begin{itemize}
\tightlist
\item
  the framework has begun to isolate the right kinds of mathematical
  objects and interfaces
\item
  it has not yet resolved the core proof obligations that would turn
  those interfaces into a proof
\end{itemize}

\end{document}
